\section{Local deformations and the quartic impurity}\label{app:interacting}
\subsection{Finite-depth and finite-time deformations}

\begin{proposition}[Preservation of the represented half-chain pair]\label{prop:factor}
Suppose $\alpha$ is an automorphism of $\A$ of the form
\[
\alpha=\Ad(u)\circ(\beta_L\otimes\beta_R),\qquad
u\in\mathcal U(\A),\quad\beta_s\in\operatorname{Aut}(\A_s).
\]
The half-chain von Neumann algebras for $\omega\circ\alpha$ are simultaneously
unitarily conjugate to $M_L,M_R$. In particular their types and the equality
$M_L=M_R'$ are preserved. Under this identification the new state vector is
$\pi(u)^*\Omega$. If $\alpha$ preserves each $\A_s$ setwise, one may take $u=1$,
and all state-specific game values and attainment properties are preserved.

The displayed factorization holds for every circuit of finite depth whose
layers consist of disjoint gates of uniformly bounded diameter. It also holds
for a fixed-time dynamics generated by an interaction whose two restrictions to
the halves generate dynamics and whose terms meeting both halves are absolutely
summable in operator norm.
\end{proposition}

\begin{proof}
The GNS triple for $\omega\circ\alpha$ can be realized as
$(\pi\circ\alpha,\HH,\Omega)$ because $\alpha$ is onto. In this representation
the two algebras are $\pi(u)M_s\pi(u)^*$. Applying $\pi(u)^*$ to the Hilbert
space proves the first claims. If $\alpha$ preserves each half algebra, its
restrictions determine $\alpha$ on the tensor product, with $u=1$.

In a circuit layer the gates wholly in each half define the two half-chain
automorphisms. All gates meeting both halves lie in a finite interval around
the cut, so their product is one local unitary. Disjoint gates commute, giving
the displayed form for each layer. Products retain the form by
$\beta\circ\Ad(v)=\Ad(\beta(v))\circ\beta$. The crossing unitary for the product
is local because a finite-depth circuit maps each local observable into another
finite interval. For the dynamics, remove the boundary sum $v$ to get the
product dynamics $\gamma_t=\gamma_{L,t}\otimes\gamma_{R,t}$. The full dynamics is
a bounded perturbation of $\gamma_t$ by $v=v^*\in\A$. Its Dyson cocycle
$u_t\in\mathcal U(\A)$ gives $\alpha_t=\Ad(u_t)\circ\gamma_t$.
\end{proof}

The factorization is related to the local automorphism factorizations of Refs.~\cite{NO,Ogata}. The explicit bulk applications are proved in Theorem~\ref{thm:bulk-circuit}.

\subsection{The quartic impurity}
\begin{proof}[Proof of Theorem~\ref{thm:impurity}]
\emph{Uniform finite-volume estimate.}
On $I_m=[-m+1,m]$ for $m\geq2$, use open boundaries and write $h_m$ for the
one-particle compression. Its eigenvalues $e_k=2\cos(k\pi/(2m+1))$ are nonzero.
In the finite Jordan--Wigner representation the unperturbed Hamiltonian
minus its lowest energy is $\sum_k |e_k| b_k^*b_k$, where $b_k$ annihilates a
particle if $e_k>0$ and a hole if $e_k<0$. The free ground vector $\Omega_m$ is
the common vacuum of the $b_k$. Denote by $K_{\lambda,m}$ this vacuum-normalized
Hamiltonian plus $\lambda V$, by $E_m$ its least eigenvalue, and by $\Psi_m$
a normalized ground vector. The local even identification of $V$ with the
same spin operator $V_{\rm s}$ is independent of $m$.

The pull-through identity gives
\[
 (K_{\lambda,m}-E_m+|e_k|)b_k\Psi_m=-\lambda[b_k,V]\Psi_m.
\]
Since the operator on the left is bounded below by $|e_k|$, this implies
$\|b_k\Psi_m\|\leq |\lambda|\|[b_k,V]\|/|e_k|$. For normalized eigenvectors
$u_k$ of $h_m$, commuting $b_k$ across the four factors of $V$ by the
anticommutation relations gives the bound
\[
 \|[b_k,V]\|\leq 2\bigl(|\langle u_k,f\rangle|+|\langle u_k,g\rangle|\bigr).
\]
Here each remaining factor is a contraction. On these volumes
$f=h_m\delta_0/\sqrt2$ and $g=h_m\delta_1/\sqrt2$, so Parseval yields
\[
 \sum_k\|b_k\Psi_m\|^2
 \leq 8\lambda^2\sum_k
 \frac{|\langle u_k,f\rangle|^2+|\langle u_k,g\rangle|^2}{|e_k|^2}
 =8\lambda^2.
\]
The complement of the vacuum projection is bounded by the excitation-number
operator. Consequently
$|\langle\Omega_m,\Psi_m\rangle|^2\geq1-8\lambda^2$. The estimate holds for
every ground vector and also implies finite-volume nondegeneracy, since an
eigenspace of dimension at least two would contain a vector orthogonal to
$\Omega_m$.

\emph{A normal ground component.}
Let $\omega_0$ be the infinite XX ground state and $\psi$ any weak-$*$ limit of
the finite perturbed ground states, considered on local spin observables.
The limit is a ground state for the perturbed dynamics. Indeed the
finite-volume ground-state inequalities pass to the limit on the local
algebra, a generator core for a finite-range spin interaction \cite{BR2}.
The free finite-volume states converge locally to $\omega_0$: the compressions
$h_m$, extended by zero, converge strongly to $h$, and the negative spectral
projections converge strongly because $h$ has no spectral atom at zero.
Wick's rule then gives convergence on the even spin local algebra, while
both states vanish on odd observables. The finite-volume overlap bound
therefore gives
\[
 \|\psi-\omega_0\|\leq 2\sqrt8\,|\lambda|<2,
\]
first on local observables and then by norm density.

Let $z\in\A^{**}$ be the central support of the XX GNS representation
$(\pi_0,\HH_0,\Omega_0)$. Since $\omega_0$ is pure, the corner $z\A^{**}$ is
normally isomorphic to $B(\HH_0)$. The normal extensions of the states
satisfy $\psi(z)>0$: if it were zero, evaluation on $2z-1$ would give
$\|\psi-\omega_0\|=2$. The projection $z$ is invariant under the perturbed
dynamics. To see this, write the unperturbed dynamics in the XX representation
as $\Ad(e^{itK_0})$. Bounded perturbation by $\lambda V_{\rm s}$ is implemented
there by $K=K_0+\lambda\pi_0(V_{\rm s})$, using the same Dyson cocycle on
$\A$. The representation after each perturbed automorphism is therefore
unitarily equivalent to $\pi_0$ and has the same central support.

The normalized component $a\mapsto\psi(za)/\psi(z)$ is a normal ground
state in $B(\HH_0)$ for $\Ad(e^{itK})$. Here is the positive-energy argument
for this step. Extend the GNS representation of $\psi$ normally to $\A^{**}$
and let $p$ be the image of $z$. The ground-state implementing unitaries
fix its cyclic vector and commute with $p$ by invariance of $z$. They thus
have nonnegative generator on $p\HH_\psi$, with invariant normalized vector
$p\Omega_\psi/\sqrt{\psi(z)}$. This is the GNS space of the component,
and the normal representation there factors through $z\A^{**}$. It is also
the GNS space of the corresponding normal state on $B(\HH_0)$, because
$\pi_0(\A)$ is strongly dense and the representation is normal.

Write this normal state as a trace-class density $D$. Its invariance implies
$D$ commutes with $e^{itK}$. Each positive-eigenvalue eigenspace of $D$ is
finite-dimensional and invariant under $e^{itK}$. It therefore has a basis
of eigenvectors of $K$. If $u$ is one such unit vector with $Ku=eu$ and
$Du=d u$, $d>0$, then in the Hilbert--Schmidt model of the component GNS
space the subspace spanned by $|v\rangle\langle u|\sqrt D$ for $v\in\HH_0$
carries the implementing group $e^{it(K-e)}$. Its generator must be
nonnegative, so $e=\inf\spec(K)$. Thus $K$ has a ground eigenvector $\Psi$.
Every such vector is cyclic for the irreducible spin representation
$\pi_0$, so its spin state is pure and has the original half-chain pair.

\emph{The interaction changes the state and leaves gapless excitations.}
The operator $V$ is a two-body term in the original fermions, with a nonzero
two-particle matrix element and zero action on the bare zero- and one-particle
spaces. It also acts nontrivially on the free Fermi sea. If $P$ is the negative
spectral projection of $h$, then $(1-P)f,(1-P)g$ are linearly independent, as
are $Pf,Pg$, since in Fourier space the two form factors are proportional to
$\cos k$ and $e^{ik}\cos k$ on each Fermi arc. The two-particle, two-hole
component of $\pi_0(V_{\rm s})\Omega_0$ is the wedge product of these two
pairs (conjugating the hole pair). It is nonzero. This computation takes place
in the common even cyclic subspace of the spin and Fock representations
under the even Jordan--Wigner identification. The even cyclic subspaces
are identified by GNS uniqueness because the states agree on the even
algebra \cite[Sec.~6.2]{KMSW}.
It follows that $\Omega_0$ cannot be a ground eigenvector of $K$ for
$\lambda\ne0$. The purity and Haag duality assertions follow from
$\pi_0$ and the XX theorem \cite{vLSW25}.

For gaplessness, let $\tau_x$ be spin translation. For any fixed local
observable $a$, $\pi_0(\tau_x(a))$ converges weakly to $\omega_0(a)1$ as
$x\to\infty$. All its weak cluster points commute with the full spin algebra
by locality, and purity and translation invariance identify the scalar.
Let $\mathcal L_0$ denote the unperturbed local commutator derivation without
the factor $i$. For any $e\in(0,2)$ and any $\eta>0$, choose a local even
$b$ with $\|\pi_0(b)\Omega_0\|\geq1/2$ and
$\|\pi_0(\mathcal L_0 b-e b)\Omega_0\|<\eta$. Such $b$ comes from
$a^*(f_+)a(g_-)$, with normalized wave packets $f_+,g_-$ supported in
spectral bands of $h$ near $e/2,-e/2$, followed by finite spatial truncation.
The free energy of this particle--hole vector is within the chosen band width
of $e$. Truncation controls the observable and the commutator in norm because
$h$ is bounded. Translate $b$ away from the impurity, so
$[V_{\rm s},\tau_x(b)]=0$. Local observables are differentiable for the
quasi-local dynamics, so their representatives preserve the domain of $K$
and satisfy the local commutator identity on it. Writing $E=\inf\spec(K)$,
one has
\[
 (K-E-e)\pi_0(\tau_x(b))\Psi
 =\pi_0\bigl(\tau_x(\mathcal L_0 b-e b)\bigr)\Psi.
\]
The squared norms of the translated vectors on both sides of this estimate
converge to the corresponding squared norms at the free vacuum, by the
weak limit for translates applied to the appropriate local products.
Taking $\eta\to0$ and normalizing gives approximate eigenvectors of $K-E$
at $e$. Hence all $e\in(0,2)$ belong to its spectrum. The game claim follows from the half-chain pair and the proof of Theorem~\ref{thm:value}.
\end{proof}

\begin{proof}[Proof of Corollary~\ref{cor:nongaussian}]
Let $Q$ be the ground spectral projection of $K$ and let $p_0$ be the
rank-one projection onto $\Omega_0$, viewed in the central corner
$z\A^{**}\cong B(\HH_0)$. The norm bound in the theorem, applied to
$2p_0-1$, gives $\psi(p_0)\geq1-\sqrt8|\lambda|$. The density of the normal
component is supported on $Q$, so $\psi(p_0)\leq\langle\Omega_0,Q\Omega_0\rangle$.
Hence $\Psi_\lambda=Q\Omega_0/\|Q\Omega_0\|$ is well defined and converges
in norm to $\Omega_0$ as $\lambda\to0$. The gauge rotations fix $\Omega_0$
and commute with $K_0$ and the density interaction, so they fix
$\Psi_\lambda$. In particular $\Psi_\lambda$ is in the even spin subspace.
In the corresponding even Fock subspace its vector state extends to the
full CAR algebra. The extension is even and gauge invariant and is pure,
because every nonzero vector is cyclic in the irreducible free Fock
representation.

Suppose this fermionic state were quasifree. Its covariance $P'$ is then a
projection by purity. On $\operatorname{span}\{f,g\}$ the free covariance
has diagonal entries $1/2$ and off-diagonal entry of modulus $4/(3\pi)$,
as follows by integrating $2\cos^2 k$ and $2e^{ik}\cos^2 k$ on the occupied
arc with measure $dk/(2\pi)$. For a normalized linear combination $r$ of
$f,g$, the expectation of $a^*(r)a(r)$ changes by at most
$2\sqrt{\sqrt8|\lambda|}$. This is the vector-state norm bound from the
squared overlap with the free vacuum. The stated range makes that change
smaller than $1/2-4/(3\pi)$. Both $P'$ and $1-P'$ are therefore strictly
positive on $\operatorname{span}\{f,g\}$, so the two particle and hole
form-factor pairs remain linearly independent.

Use the particle--hole Fock realization of this hypothetical pure
quasifree state and denote its vacuum by $\Phi$. The operator
\[
 B=a^*((1-P')f)a^*((1-P')g)a(P'f)a(P'g)
\]
is a product of four excitation creators, and $B\Phi\ne0$. The four-creation
component of $V\Phi$ is $\pm B\Phi$. Therefore the vacuum expectation of
$[V,B^*]$ is $\mp\|B\Phi\|^2\ne0$. The vacuum expectation of
$\mathcal L_0(B^*)$ is zero because the quadratic derivation changes only one
annihilation factor at a time to a linear combination of excitation
creators and annihilators, while the other three factors still annihilate
excitations. Every such term changes excitation number and has zero vacuum
expectation. The even CAR observable $B^*$ is in the domain of the spin
commutator derivation after Jordan--Wigner. Finite spatial truncations of
its smearing vectors approximate both $B^*$ and its unperturbed commutator
in norm since $h$ is bounded, and the impurity commutator is bounded.
Stationarity of the ground state requires the expectation of
$\mathcal L_0(B^*)+\lambda[V,B^*]$ to vanish, a contradiction.
\end{proof}

\begin{proof}[Proof of Corollary~\ref{cor:impurity-finite}]
Use a common ambient Fock space over
$\mathfrak k_* = \ell^2(\mathbb Z)\oplus\overline{\ell^2(\mathbb Z)}$.
Let $P_m^\pm$ be the projections onto the positive and negative spectral
subspaces of the open $h_m$, extended by zero outside $I_m$. The finite
quasiparticle space is
$\mathfrak k_m=P_m^+\ell^2\oplus\overline{P_m^-\ell^2}$.
For the spin finite-volume ground vectors $\Psi_m$ use the finite
Jordan--Wigner identification with their quasiparticle Fock vectors and
embed these into the ambient Fock space. The free finite vacua all map to
the same ambient vacuum. The infinite free particle--hole space is
$\mathfrak k=(1-P)\ell^2\oplus\overline{P\ell^2}$. On fixed local even
spin observables the finite Fock realizations converge in operator norm
to the infinite one when viewed as polynomials in ambient annihilators
and creators, since $P_m^\pm$ converge strongly. The finite vectors are
even in excitation parity. The Hamiltonian commutes with parity, its
ground vector is unique by the overlap bound, and the overlap with the
even free vacuum is nonzero.

We show precompactness of these vectors. Write $v_k$ for the excitation
one-particle basis, $\epsilon_k=|e_k|$, and $T_m=K_{\lambda,m}-E_m\geq0$.
The CAR commutator used above has the form
\[
 [b_k,V]=\sum_{\ell=1}^4\langle v_k,w_{m,\ell}\rangle A_{m,\ell},
 \qquad \|A_{m,\ell}\|\leq 1,
\]
where the $w_{m,\ell}$ are the particle and hole components of $f,g$
in $\mathfrak k_m$. Each $A_{m,\ell}$ is the remaining signed product of
three operators and is independent of $k$. Thus the vector of annihilations
in $\mathfrak k_m\otimes\Gamma(\mathfrak k_m)$ is
\[
 \sum_k v_k\otimes b_k\Psi_m
 =-\lambda\sum_{k,\ell} v_k\langle v_k,w_{m,\ell}\rangle\otimes
     (T_m+\epsilon_k)^{-1}A_{m,\ell}\Psi_m .
\]
For $\epsilon_k<\delta$ its squared norm is bounded by
$8\lambda^2$ times
\[
 \frac12\sum_{j=0}^1
 \big\langle\delta_j,\mathbf1_{(-\delta,\delta)}(h_m)\delta_j\big\rangle.
\]
This is the same estimate as for the total excitation number, restricted
to these energies. The displayed bound tends to zero as $\delta\downarrow0$
after taking $\limsup_m$, by strong convergence of these spectral
projections and absence of a free zero-energy atom.

For $\epsilon_k\geq\delta$, partition $[\delta,2]$ into finitely many energy
intervals $J$ of diameter at most $\eta$, and replace the resolvent by
$(T_m+c_J)^{-1}$ with $c_J\in J$. Its norm error is at most $\eta/\delta^2$
by the resolvent identity. The resulting vector lies in
$F_m\otimes\Gamma(\mathfrak k_m)$, where $F_m$ is the span of all energy-band
projections of the four $w_{m,\ell}$. Its error in norm on the high-energy
part is at most $4|\lambda|\eta/\delta^2$, using the orthogonality of the
$v_k$ and $\|w_{m,\ell}\|\leq1$. With the low-energy bound this gives
arbitrarily small $\limsup_m$ distance to $F_m\otimes\Gamma(\mathfrak k_m)$.

In Fock space the expectation of the number operator outside $F_m$ is
the squared norm of the part of the annihilation vector orthogonal to
$F_m$ in the first tensor factor. That number operator bounds the
complement of the Fock projection onto $\Gamma(F_m)$. Hence $\Psi_m$
has arbitrarily small $\limsup_m$ distance to $\Gamma(F_m)$.
The finitely many vectors spanning $F_m$ converge in the ambient
one-particle space to the corresponding infinite-volume band projections.
They are linearly independent in the limit: different bands and particle
and hole sectors are orthogonal, and within a band the two Fourier
form factors $\cos k$ and $e^{ik}\cos k$ are independent on a set of
positive measure. The partition can be chosen with each interval of
positive length, and the free dispersion has no atoms at the endpoints.
Consequently their span projections converge in norm to a finite-rank
projection onto $F\subset\mathfrak k$. The projections onto the finite
Fock spaces $\Gamma(F_m)$ also converge in norm to that onto $\Gamma(F)$,
as can be seen by choosing convergent orthonormal bases and taking wedge
products. For any tolerance, the vectors thus lie within that tolerance
of a fixed finite-dimensional subspace for all sufficiently large $m$.
This proves norm precompactness, with every limit vector a unit even
vector in $\Gamma(\mathfrak k)$.

Each such limit determines the spin limit state using the common even
GNS subspace and vanishing odd expectations. It is a pure vector state
in the XX representation and a ground state by the local ground-state
inequalities. If the finite-chain game value did not tend to one, a
subsequence with a uniform positive loss floor would have a further
subsequence converging locally to one of these limit states. That state
has value one by its half-chain pair. By local POVM approximation as in
Lemma~\ref{lem:local}, it admits a finite-support strategy with loss smaller
than the putative floor. Its loss in the convergent finite states tends
to that same value, a contradiction. The same argument with local embezzling unitaries, as in the finite-chain proof of Theorem~\ref{thm:value}, proves convergence of the embezzlement error.
\end{proof}
